\section{R16: a source-pinned native word machine}
\label{sec:r16-wqk}

R16 imports the finite TOMAGI 1.0 execution semantics from Tom Klootwijk's
World Query Kernel 0.6 package and gives them a native CUDA realization.
Its Cell48 operators execute against resident State64 words. An explicit
fresh-emission binding connects that machine to the inherited two ASA/NA
mask stages and JK state; a separate exact rational adapter discovers
affine plane crossings. These are two bounded engine profiles within the
retained aTOMos formalization.

The source authorities are \path{src/c/tomagi.h},
\path{src/c/tomagi.c}, the source Python TOMAGI implementation, and
\path{spec/TOMAGI_1_0_FORMAL_DEFINITION.md} in the supplied WQK package.
The release provenance manifest pins file contents and preserves the
source attribution. No license grant or independent authorship finding
is inferred from a notice. The complete integration contract is
\path{formal/WQK_INTEGRATION.md}.

\begin{panel}
The source opcode \code{PHI} advances a finite phase. It is distinct
from the exact $a+b\varphi$ algebra. \code{SDF0} assigns zero; it does
not calculate Euclidean distance. \code{KLEIN} retains a reduced chart
and wrap parity; it does not retain the full winding lift. These precise
meanings preserve compatibility without changing the inherited physical
or geometric contracts.
\end{panel}

\subsection{Canonical words, address, and admission}

For an integer $z$, define unsigned reduction $u_{32}(z)=z\bmod2^{32}$
and signed interpretation $s_{32}$ of that word. Thus
\begin{equation}
 \wrap_{32}(z)=s_{32}(u_{32}(z)),\qquad
 \operatorname{mod}_m(z)=z-m\lfloor z/m\rfloor,\quad m>0.
 \label{eq:r16-word-wrap}
\end{equation}
Modulo addition of unsigned raw words implements the source wrapping
addition without signed overflow. Where the source specifies a wider
intermediate, reduction occurs after that intermediate is evaluated.
Signed division in \code{LSYS} truncates toward zero; phase and tick
wrap counts use floor division. They are different operations for
negative inputs.

The canonical little-endian stream contains a 128-byte header and $N$
48-byte cells, exactly $128+48N$ bytes. It declares the eight-byte magic
\code{TOMAGI1\textbackslash0}, version \code{0x00010000}, cell size 48,
state size 64, and six zero reserved words. Admission requires $N>0$,
an entry below $N$, opcodes in $[0,15]$, successors below $N$, and
strictly increasing unique unsigned keys $(K_{hi},K_{lo})$.

State64 comprises sixteen raw 32-bit words. The signed coordinate and rate
fields are
\[
 (\rho,\theta,t,\phi,\;v_\rho,v_\theta,v_t,v_\phi).
\]
They are followed by unsigned orientation, sheet, branch, cell, lineage, and output;
residual is signed and status unsigned. Cell48 has twelve words:
two key words, opcode, flags, four signed arguments, two successors,
payload, and auxiliary word.

The constructor's \code{header\_state()} retains the entire initial
state. Its \code{entry\_state()} changes only the cell index to the
declared entry; that is the constructor's initial lane. Explicit
\code{set\_states()} preserves supplied raw words and resets sidecar
faults, receipts, and dispatch epoch. It does not normalize an initial
branch or phase before the first instruction.

With periods
\begin{equation}
 (M_\rho,M_\theta,M_t,M_\phi)=(2^{20},2^{18},2^{14},2^{12}),
 \label{eq:r16-periods}
\end{equation}
the pre-instruction key is the exact bit-field concatenation
\begin{equation}
 \begin{split}
 K={}&\operatorname{mod}_{M_\rho}(\rho)2^{44}
      +\operatorname{mod}_{M_\theta}(\theta)2^{26}\\
    &+\operatorname{mod}_{M_t}(t)2^{12}
      +\operatorname{mod}_{M_\phi}(\phi).
 \end{split}
 \label{eq:r16-address}
\end{equation}
It is a finite address, not a metric. Physical coordinates, units, and
chart identifications need an explicit application binding.

\subsection{Instruction bodies and the mandatory postlude}

Each row below denotes the instruction body before the shared postlude.
The source write order is part of its semantics: \code{SET} and
\code{JIT1} can address any state field, including branch, status,
cell, and lineage. A JIT1 target that aliases branch can modify its
earlier parity assignment.

{\small
\begin{longtable}{@{}P{.19\textwidth}P{.77\textwidth}@{}}
\toprule
Opcode & Source body retained in R16 \\
\midrule
\endfirsthead
\toprule Opcode & Source body retained in R16 \\
\midrule
\endhead
0 NOP & No body change. \\
1 SET & Write argument bits to field \code{flags\,\&\,15}. \\
2 JIT1 & Source hash/parity from seed and pre-state; assign branch, then add the selected signed perturbation with word wrapping. \\
3 KIN2 & Update each rate, then its position by the updated rate, modulo $2^{32}$. \\
4 PHI & Wide signed phase addition; Euclidean reduction by $2^{12}$, wrap status, selected branch and orientation action. \\
5 TIME & Wide signed tick addition; reduction by $2^{14}$, wrap-parity branch and conditional lineage mix. \\
6 SDF0 & Residual zero, ZERO status set, branch one. \\
7 CONE & Radial/angular interval maximum; wrap residual first, then test signed residual $\leq0$. \\
8 SPHERE & Radial and optional cyclic-phase interval maximum; wrap residual first, then test signed residual $\leq0$. \\
9 KLEIN & Reduce signed rho; odd wrap parity applies the selected theta map, phase negation, orientation and optional sheet flip. \\
10 RADIX & Select the declared bit of the pre-instruction key; shifts outside $[0,63]$ refuse. \\
11 HINGE & If branch low bit is one, add coordinate arguments with wrapping and apply selected flips. \\
12 LSYS & Signed orientation/branch phase turn and rate division by $2^{\operatorname{clamp}(arg_1,0,30)}$, truncating toward zero. \\
13 PROJECT & Assign payload to output. \\
14 EMIT & Assign payload, set EMIT, optionally HALT. \\
15 HALT & Set HALT. \\
\bottomrule
\end{longtable}}

The \code{CONE}/\code{SPHERE} residual wraps before its inside test.
R16 preserves that finite behavior. An application wanting the
unwrapped interval geometry must additionally prove a no-overflow
domain. Neither opcode is a general Euclidean signed-distance field.

For \code{PHI}, for example, the exact wider intermediate is
\begin{equation}
 r=s_{32}(\phi)+arg_0,\qquad
 w=\lfloor r/2^{12}\rfloor,\qquad \phi_{\rm body}=r-w2^{12}.
 \label{eq:r16-phase}
\end{equation}
The branch is the selected wrap-parity or half-phase bit. This is not
a multiplication by the golden ratio. Likewise, \code{KLEIN} uses the
parity of $\lfloor s_{32}(\rho)/2^{20}\rfloor$ and does not preserve
that integer after reduction. A 32-bit lineage hash cannot recover
an unbounded winding count.

Every successful body, including HALT and EMIT-HALT, then normalizes
$\theta,t,\phi$, masks orientation and branch to one bit, and updates
lineage. It does not globally normalize rho or sheet. Writing $i^-$
for the executed cell and $K^-$ for the pre-instruction key,
\begin{equation}
 \begin{split}
 L^+=\operatorname{mix32}\bigl(&L_{\rm body}\oplus payload\oplus aux
       \oplus K_{hi}^-\\
       &\oplus\operatorname{rotl}_{32}(K_{lo}^-,7)
       \oplus b^+\oplus i^-\bigr).
 \end{split}
 \label{eq:r16-lineage}
\end{equation}
The source unsigned hash retains its exact constants and operations.
If the resulting state is not halted, the normalized branch chooses
a successor. With REKEY enabled, an exact lookup of the new key takes
precedence on a hit; a miss sets REKEY\_MISS and uses that successor.
HALT takes no successor. An already halted lane executes no instruction.

The native fault sidecar reports bad cell, opcode, or RADIX shift. A
refused transition publishes an error receipt and leaves State64
unchanged. This result is separate from a successful source transition.
The reviewed source OpenCL HINGE uses a whole-word nonzero branch test,
where C/Python use \code{branch\,\&\,1}. Raw initial branch $2$ exposes
the difference. R16 follows the C/Python low-bit rule; agreement on
already-normalized states would not detect this discrepancy.

\subsection{Literal one-bit planes and immutable texture records}

For 32 source words $w_j$, with zero padding only for missing terminal
lanes, define the transport and its inverse:
\begin{equation}
 P_b=\sum_{j=0}^{31}\bigl((w_j\gg b)\mathbin\&1\bigr)2^j,
 \qquad
 w_j=\sum_{b=0}^{31}\bigl((P_b\gg j)\mathbin\&1\bigr)2^b.
 \label{eq:r16-planes}
\end{equation}
The two transformations are inverse because both address the same bit
$(j,b)$. Signs, flags, operands, and instructions are preserved. This
packing changes storage order, not operator meaning or arithmetic precision.

The native upload reconstructs the header and cells on the GPU from
raw integer plane textures. Decoded Cell48 records remain immutable,
banked within device limits, and are read by three
\code{tex1Dfetch<uint4>} operations per cell. The matching global-memory
variant reads the same records. Element reads use integer indices and
no filtering or normalized-value conversion. The canonical source,
decoded program, and mutable State64/receipt arrays are separate objects.

Program-bank capacities obey device resource limits and the signed
texture-index range. Integer texture access establishes use of that
hardware path, not a cache-hit rate, permanent cache residency, or
a speed guarantee. Memory access and execution measurements belong
to the recorded evidence. This profile does not extend the lowering
domain of every inherited exact expression or transcendental operator.

\subsection{Fresh EMIT and actual state feedback}

The versioned binding \code{ATOMOS-TOMAGI-EMIT-FEEDBACK-R1} wraps the
pure VM. Before an enabled coupled step, it copies the low and high
32-bit halves of the old hinge word $q_k$ to raw State64 rho and
$v_\rho$, respectively; halted or faulted lanes are skipped. This is
an explicit application state mapping, not an implicit change to the
pure TOMAGI transition. The VM can still run without the binding.

After one canonical instruction, a fresh receipt is eligible exactly
when that instruction executed successfully and was EMIT:
\begin{equation}
 e_{k,i}=[executed_{k,i}=1]\land[opcode_{k,i}=14]
                  \land[error_{k,i}=0].
 \label{eq:r16-emission}
\end{equation}
Each $(\text{epoch},\text{lane})$ is consumed at most once. The drive
word is the zero-extended 32-bit payload. PROJECT, an instruction after
an earlier EMIT, an already halted lane, a refusal, or a duplicate
receipt supplies no fresh event. EMIT-HALT supplies exactly one.
Equal payloads in different successful EMIT epochs remain distinct events.

Let $P$ be the present mask, $M=2^{64}-1$, and let the declared Boolean
word body give $x=X(q_k,d_k)$. The literal two mask stages are
\begin{align}
 a_i(x)&=x\mathbin\&A_i\mathbin\&P,
 &n_i(x)&=a_i(x)\mathbin\&N_i,\nonumber\\
 T_i(x)&=\begin{cases}n_i(x),&n_i(x)\mathbin\&B_i=0,\\
                         0,&\text{otherwise},\end{cases}
 &y_k&=T_1(T_0(x)).
 \label{eq:r16-mask-stages}
\end{align}
Here NA is the retained mask stage, not logical complement. The
declared $J,K$ word bodies read the same old $q_k$ and the derived
$y_k$, then commit
\begin{equation}
 q_{k+1}=\bigl[(J_k\mathbin\&(M\oplus q_k))\mathbin\vert
      ((M\oplus K_k)\mathbin\&q_k)\bigr]\mathbin\&P.
 \label{eq:r16-emit-jk}
\end{equation}
An accepted emission toggles the separate event parity once. This
fixed-chart binding leaves chart orientation zero; VM orientation
remains its own state field and is not a geometric seam declaration.

Ordered native dispatches make the committed $q_{k+1}$ available to
the next enabled pre-step binding, without a whole-state host readback.
The coupler's lifetime, state-generation, and expected-epoch checks
detect replacement or advancement through the managed VM API. Direct
DeviceView writes or low-level launches remain caller-owned operations
and must not bypass this ownership while coupled execution is active.
Dispatches sharing mutable lanes require one ordered stream or explicit
cross-stream dependencies; host call order alone does not order distinct
CUDA streams. These conditions preserve the declared execution sequence;
they do not make the finite recurrence a physical measurement.

\subsection{What is deliberately outside this import}

The WQK finite learner's exact candidate survival, explicit ambiguity,
and separate validation gates are useful design material. Its delivered
candidate families remain evaluated source material in R16. They are
not imported as additional CPU toy families or claimed as native GPU
model search. The source compiler's \code{formal.evaluate} can compute
an answer in Python and lower its serialized bytes to EMIT cells.
Running those cells executes the materialization, not the preceding
Python search or proof.

A future exact calibration profile would need to retain unresolved
eligible candidates, not discard overflow or unsupported evaluation as
a mismatch. Unique selection requires exactly one verified fitting
candidate and none unresolved. Noise, observability, physical units,
and independent observations remain part of that application's model.
The retained physics chapters supply their existing physical equations;
this source import supplies a concrete finite execution and event interface.
